\paragraph{What the claim covers.} Every cost result here concerns
\textbf{centralised-inference computational cost}: the wall-clock latency and
memory of executing the communication module on a single device, where all
agents' observations are available to one process. It is \emph{not} a claim
about bandwidth-constrained distributed deployment. In that regime the
objective is different --- reducing the number of messages physically
transmitted over a link --- and top-$k$ sparsification may well be the correct
design even though it does not reduce local compute. Conflating the two would
misrepresent this result, and the analytic ceiling of
Proposition~\ref{prop:ceiling} says nothing about transmitted bytes.

\paragraph{Measurement envelope.} Single-round attention; scaled dot-product
attention with the fused backends available in PyTorch 2.11 / CUDA 12.8; input
embedding width $64$; head width $d \in \{16,32,64,128,256\}$; $N \le 512$;
batch $32$--$4096$; \texttt{fp32} and \texttt{bf16}; eager, \texttt{torch.compile}
and CUDA-graph execution. Hardware: one GPU architecture (NVIDIA Blackwell,
sm\_120) and one x86\_64 CPU. We make no claim outside this envelope.

\paragraph{No cross-generation GPU replication.} The strongest available
robustness evidence is the CPU result, whose machine balance differs from the
GPU's by more than an order of magnitude, plus an idle second device of the
same model. We were unable to replicate on a different GPU generation --- the
older hardware that produced this project's earlier runs is no longer
available. ``Structurally'' versus ``on this architecture'' is therefore
supported by the roofline argument and the CPU agreement rather than by a
second GPU generation, and a reader who weights that evidence differently
should discount the generality accordingly.

\paragraph{Unlocked clocks.} We lacked the privilege to lock GPU clocks. We
compensate with interleaved timing, minimum-over-samples reporting, published
dispersion (Appendix~\ref{app:dispersion}) and an explicit drift control
(Section~\ref{sec:drift}), and we report that absolute latency under
concurrent load can move by tens of percent. Ratios are far more stable, and
all claims are stated in ratios.

\paragraph{Statistical power of the training study.} Five seeds per arm gives a
minimum attainable two-sided permutation $p$ of $0.0079$ against a Holm
threshold of $0.0083$ for six contrasts. This is enough to detect only perfect
separation, and it is a consequence of the compute budget, not a property of
the methods. The six null contrasts of Section~\ref{sec:training} should be
read as ``this design could not resolve an effect of this size'', and no
equivalence between sparse and dense is claimed anywhere. A study powered to
resolve the $d \approx 0.6$ effect suggested by C1 would need roughly $45$
seeds per arm.

\paragraph{Environment.} The partially observable \texttt{simple\_spread}
variant is our own, documented in Section~\ref{sec:training} and released. It
carries a difficulty confound across $N$ that we measure and publish, and which
is why no across-$N$ return comparison appears in this paper. Results on MPE
tasks may not transfer to environments with different observation structure.

\paragraph{Implementation scope.} \Gath{} is a straightforward PyTorch
implementation using \texttt{gather} and a batched contraction; it is not a
hand-written fused kernel. A custom kernel that fuses selection, gather and
contraction would reduce the memory traffic that
Section~\ref{sec:roofline} identifies as the binding constraint, and could
plausibly close part of the gap. It could not, however, escape
Proposition~\ref{prop:ceiling}, nor avoid ranking all $N$ scores. We regard
such a kernel as the most informative single follow-up to this work, together
with the structural-sparsity direction of Proposition~\ref{prop:escape}.
